% !TEX root = tkf.tex
% Results subsection: empirical evidence of the structural bias
% in the BP cumulant under column-factorised q. Theory derivation
% lives in Appendix~\ref{app:cumulant-structural-bias}.

\subsubsection{Structural bias of the BP cumulant under
column-factorised \texorpdfstring{$q$}{q}: empirical results}
\label{sec:cumulant-structural-bias-results}

The SVI-VarAnc refinement of cherry-trained parameters
(Section~\ref{sec:results-svi-varanc}) systematically inflates the
TKF91 / TKF92 indel rates by a factor of $\beta_\lambda \approx
2.85$, $\beta_\mu \approx 2.32$ relative to truth on simulated
data --- a property of the column-factorised variational family
itself, not finite-data noise.  Appendix~\ref{app:cumulant-structural-bias}
gives the precise derivation: the BP cumulant tensor
$\mathbf{W}_e$ is an exact computation of $\expect_q[\#\text{chain
transitions}]$ under factorised $q$, but the truth posterior
$p^*(\boldsymbol{\zeta} \mid \text{data})$ does not factorise across
columns (because the per-branch WFST chain has positive
$T_e[I, I]$ and $T_e[D, D]$ entries even at $\ext = 0$), so the
factorised cumulant differs from $\expect_{p^*}[\#\text{chain
transitions} \mid \text{data}]$ by an integrated connected-correlation
term (\eqref{eq:cumulant-bias-correct} in the appendix).  This
subsection presents two diagnostics that quantify the bias on a
TKF91 simulation.

\paragraph{Diagnostic 1: bias persists at $q$ = exact truth posterior.}
Setting $q$'s edge conditionals $\mathbf{q}_{\text{cond}}^{(e)}$ equal
to the data-independent TKF91 prior conditional matrix
$\mathbf{T}_{e,\text{prior}}$ and running BP gives per-column
pair-marginals exactly equal to the per-column Felsenstein truth
posterior.  Even at this $q$, the cumulant
$\mathbf{W}$ (eq.~\ref{eq:cumulant-W}) has a substantial bias against the
labelled-simulation ("oracle") chain count.  A diagnostic
sweeps $n_{\text{fams}}
\in \{1, 5, 10, 30\}$ on a 6-leaf caterpillar tree (branches
$0.3$, $\lambda = 0.04$, $\mu = 0.05$, $\ext = 0$, root
length $L = 240$):

\begin{table}[ht]
\centering\small
\begin{tabular}{rrrrrr}
\toprule
$n_{\text{fams}}$ & oracle total & $\|\mathbf{W} - \mathbf{O}\|_1$ & rel.\ &
$W[I, I]$ (oracle) & $W[D, D]$ (oracle) \\
\midrule
1 & 2470 & 117.6 & 4.76\% & 20.7 (0)  & 4.5 (2)   \\
5 & 12{,}113 & 521 & 4.30\% & 85.5 (0) & 12.9 (8) \\
10 & 24{,}461 & 1120 & 4.58\% & 189.6 (4) & 20.5 (10) \\
30 & 72{,}685 & 3721 & 5.12\% & 629.3 (14) & 82.2 (22) \\
\bottomrule
\end{tabular}
\caption{Cumulant total $\mathbf{W}$ vs labelled-simulation oracle
$\mathbf{O}$ at $q_{\text{cond}}^{(e)} =$ TKF91 prior conditional
(i.e., $q$ pair-marginals exactly equal the per-column Felsenstein
truth posterior).  Bias scales linearly with $n_{\text{fams}}$
($31.65\times$ growth for $30\times$ growth in $n_{\text{fams}}$,
versus $\sqrt{30} = 5.48\times$ if it were Monte-Carlo noise),
confirming the bias is structural --- linear scaling reflects the
sum of independent per-family biases, each of which has the same
non-zero expected value because the per-family bias is itself a
sum of $\binom{L}{2}$ pair-correlation terms.  The $W[I, I]$ entry
is $45\times$ over-counted at $n_{\text{fams}} = 30$ relative to
the oracle.  Common entries $W[M, M]$ and $W[S, M]$ match the
oracle within $1.5\%$ (not shown).}
\label{tab:cumulant-structural-bias}
\end{table}

\paragraph{Diagnostic 2: bias does not shrink with Adam convergence.}
A separate diagnostic sweeps Adam inner
iterations on the variational $q$ with untied per-(edge, column)
logits and confirms the cumulant bias plateaus at the same $\sim
4\%$ relative level by $n_{\text{Adam}} = 300$ and does not shrink
with further iteration (we tested up to $n_{\text{Adam}} = 1000$).
The variational fit is not the bottleneck: the cumulant
approximation is.

\paragraph{Diagnostic 3: Fitch parsimony + BDI MLE recovers truth.}
To localise the bias to the variational machinery (and not the
downstream BDI MLE), we run the simplest possible alternative:
hard ancestral assignment by Fitch parsimony, then BDI MLE on the
resulting fully-observed WFST chain.  Specifically, for each
column we compute the per-internal-node presence/absence using
standard Fitch parsimony (postorder: state-set is the
intersection of child sets if non-empty, else union; preorder:
each node's state is the parent's state if compatible with the
child set, else the smaller compatible state).  This gives a
fully-labelled per-branch WFST trajectory $S \to \cdots \to E$
for each column, and we count $(s, s')$ transitions directly ---
no factorised $q$, no cumulant.  We then apply the same BDI M-step
used by SVI-VarAnc.

\begin{table}[ht]
\centering\small
\begin{tabular}{lrrr}
\toprule
Method & $\hat\lambda$ ($/\lambda_*$) & $\hat\mu$ ($/\mu_*$) & $\hat\ext$ \\
\midrule
Standard Fitch + BDI MLE       & $0.0508$ ($1.27\times$) & $0.0510$ ($1.02\times$) & $0.00001$ \\
Oracle (true labels) + BDI MLE & $0.0514$ ($1.29\times$) & $0.0516$ ($1.03\times$) & $0.00001$ \\
Variational SVI-VarAnc EM      & $0.114$  ($2.85\times$) & $0.116$  ($2.32\times$) & $0.054$   \\
Dollo-floor parsimony + BDI MLE & $0.103$  ($2.59\times$) & $0.104$  ($2.08\times$) & $0.00001$ \\
\bottomrule
\end{tabular}
\caption{Rate-recovery on the same TKF91 simulation
($\lambda_* = 0.04$, $\mu_* = 0.05$, $\ext_* = 0$,
$n_{\text{fams}} = 30$, $L = 240$).  Standard Fitch parsimony
combined with the same BDI MLE used by SVI-VarAnc is
\emph{indistinguishable} from the oracle at the true labels and
is $\approx 2.8\times$ closer to truth than the variational EM
fixed point.  This localises the variational bias to the
column-factorised $q$ cumulant approximation and rules out the
BDI MLE itself as the source.  The Dollo variant
(intersection-postorder $+$ parent-floods-down preorder, used
elsewhere in this paper as a "Fitch floor" diagnostic) forbids
deletion and re-attributes all rare events to insertion; this
gives the same magnitude of bias as the variational EM but for a
different mechanism (single-direction parsimony assumption rather
than cumulant approximation).}
\label{tab:fitch-plus-bdi}
\end{table}

The conclusion is unambiguous: the variational SVI-VarAnc rate
estimate is biased by the cumulant approximation, and a simpler
hard-assignment baseline (Fitch $+$ BDI MLE) is substantially
less biased on the same data.  We do not advocate Fitch $+$ BDI
as a replacement for full Bayesian inference (it discards
posterior uncertainty over ancestral states), but the result is
striking enough that we report it as the principal practical
recommendation for users who need indel-rate \emph{point
estimates} from a tree without committing to a fully Bayesian
pipeline.

\paragraph{EM converges to a biased fixed point.}
EM around the biased cumulant reaches a finite fixed point
$\theta_\infty$ rather than diverging
(Appendix~\ref{app:cumulant-structural-bias} gives a Brouwer
existence argument from the uniform bound $W_e \leq L+1$
combined with the $\theta$-independence of the BDI exposure
$T_{\text{eff}} = \sum_e t_e$).  Empirically, on the same TKF91
simulation:

\begin{table}[ht]
\centering\small
\begin{tabular}{rccc}
\toprule
Iter & $\lambda$ & $\mu$ & $\ext$ \\
\midrule
0 (truth init) & $0.040$ & $0.050$ & $0.000$ \\
1               & $0.114$ & $0.116$ & $0.0001$ \\
5               & $0.116$ & $0.118$ & $0.0004$ \\
10              & $0.116$ & $0.117$ & $0.002$ \\
20              & $0.112$ & $0.113$ & $0.015$ \\
30              & $0.107$ & $0.108$ & $0.054$ \\
\bottomrule
\end{tabular}
\caption{Tree-VBEM EM trajectory from truth-init on simulated
TKF91 data ($\lambda_* = 0.04$, $\mu_* = 0.05$, $\ext_* = 0$;
6-leaf caterpillar tree, branches $0.3$, root length $L = 240$,
$n_{\text{fams}} = 30$).  EM converges to a biased fixed point
$\theta_\infty$ within one iteration; the late-iteration coupled
drift (rates declining, $\ext$ rising) is the M-step relabelling
the spurious $I\!\to\!I$ and $D\!\to\!D$ counts as
fragment-extension events --- the only TKF92 parameter that can
absorb such repeated rare-state structure.  EM does not diverge.
We do not claim a quantitative link between the rate inflation
factor $\beta_\lambda \approx 2.85$ and the $45\times$ count
inflation on $W[I, I]$
(Table~\ref{tab:cumulant-structural-bias}); the BDI rate-MLE is a
nonlinear quotient of multiple cumulant entries and the
relationship is not first-order.}
\label{tab:em-convergence-from-truth-inline}
\end{table}

\paragraph{What this means for the rest of the paper.}
The structural bias affects only inference paths that use the
cumulant tensor as a sufficient statistic under factorised $q$ ---
specifically, the SVI-VarAnc refinement reported in
Section~\ref{sec:results-cherry-vs-msa}.  It does \emph{not}
affect (i) the ELBO value itself (a valid $\log p(\text{data})$
lower bound regardless of cumulant accuracy), (ii) ancestral
indel-presence reconstruction F1
(Section~\ref{sec:results-varanc-vs-fitch}: VarAnc and Fitch
parsimony are within $\pm 0.005$ F1 across all three difficulty
strata of the unified Pfam test), or (iii) MSA-conditioned rate
estimation via the exact-EM Maraschino training of
Section~\ref{sec:results-pfam-cherry}.  We report uncorrected
$\theta_\infty$ values throughout; the structural bias is the
principal known limitation of variational training in this paper.
